\documentclass[11pt]{article}
\usepackage[a4paper,margin=2cm]{geometry}
\usepackage{amsmath,amssymb,amsthm}
\usepackage[T1]{fontenc}
\usepackage[utf8]{inputenc}
\pagestyle{empty}
\newtheorem{theorem}{Theorem}
\newtheorem{corollary}[theorem]{Corollary}
\newtheorem{example}[theorem]{Example}
\theoremstyle{definition}
\newtheorem{definition}[theorem]{Definition}
\newcommand{\Nor}{\mathbf{Nor}}
\newcommand{\lo}[1]{\ell_1^0(#1)}
\newcommand{\ellone}[1]{\ell_1(#1)}
\newcommand{\norm}[1]{\lVert #1\rVert}
\newcommand{\BM}{d_{BM}}
\begin{document}
\begin{center}
  {\large\bfseries Extremally projective normed spaces}\\[2pt]
  Norbert Nemesh \quad $\cdot$ \quad handout for the talk
\end{center}

\paragraph{Setting.} $\Nor$ is the category of normed spaces and bounded
operators. An operator $\xi\colon X\to Y$ is \emph{$c$-topologically
surjective} if $c\,\xi(B_X^\circ)\supset B_Y^\circ$. For a set $\Lambda$,
$\ellone\Lambda$ is the Banach space of absolutely summable families and
$\lo\Lambda$ its dense subspace of finitely supported families; the spaces
$\lo\Lambda$ are the free objects of $\Nor$. $\BM$ denotes the Banach--Mazur
distance, $\overline P$ the completion of $P$.

\begin{definition}
A normed space $P$ is \emph{$C$-topologically projective} if for every
$c$-topologically surjective $\xi\colon X\to Y$ between normed spaces and every
bounded $\varphi\colon P\to Y$ there is a bounded $\psi\colon P\to X$ with
$\xi\psi=\varphi$ and $\norm\psi\le cC\norm\varphi$. It is \emph{topologically
projective} if this holds for some $C\ge1$, and \emph{extremally projective}
if it holds for every $C>1$.
\end{definition}

Extremal projectivity is Helemskii's extreme projectivity for normed spaces.
Known before: $\lo\Lambda$ is extremally projective; $P$ is extremally
projective iff it is a near-retract of some $\lo\Lambda$; then
$\overline P\cong\ellone\Gamma$ isometrically. Helemskii (2013) and
Gr{\o}nb{\ae}k (2016, Conjecture 3.1) asked whether every extremally
projective normed space is isometrically isomorphic to some $\lo\Lambda$.

\begin{theorem}[dimension restrictions]
Let $P$ be topologically projective and $Z\subset P$ a subspace. Then
$\dim Z\le\max(\operatorname{dens}Z,\aleph_0)$. Consequently $P$ contains no
infinite dimensional subspace which is complete with respect to a norm
dominating the given one; in particular no infinite dimensional Banach space
is topologically projective in $\Nor$.
\end{theorem}

\begin{theorem}[countable dimension]
Let $D\subset\ell_1$ be a dense subspace of countable algebraic dimension and
$\varepsilon>0$. There is a linear bijection $T$ of $\ell_1$ with
$\norm T\norm{T^{-1}}\le1+\varepsilon$ and $T(D)=\lo{\mathbb N}$. Hence, for a
normed space $P$ of countable dimension: $P$ is extremally projective iff
$\overline P\cong\ell_1$ or $\ell_1^n$ isometrically iff $\BM(P,\lo\Gamma)=1$
for some countable $\Gamma$.
\end{theorem}

\begin{corollary}[counterexample]
Let $v=(2^{-n})_{n\in\mathbb N}$. The space $\lo{\mathbb N}+\mathbb C v$ is
extremally projective but not isometrically isomorphic to any $\lo\Lambda$:
for $\lo S\subset E\subset\ellone S$ the extreme points of $B_E$ are the
vectors $\lambda e_s$, $|\lambda|=1$, whose convex hull is $B_{\lo S}$, so an
isometry $\lo\Lambda\to E$ forces $E=\lo S$.
\end{corollary}

\begin{example}[ladder space]
For every limit ordinal $\alpha<\omega_1$ fix $\alpha_n\nearrow\alpha$ and put
$v_\alpha=\sum_n2^{-n}e_{\alpha_n}\in\ellone{\omega_1}$,
$P_L=\lo{\omega_1}+\operatorname{span}\{v_\alpha\}$. Then $P_L$ is dense in
$\ellone{\omega_1}$, every countable piece $P_L\cap\ellone\delta$ is extremally
projective, every separable subspace of $P_L$ has countable dimension, and
$P_L$ is not topologically projective (Fodor's lemma).
\end{example}

\begin{theorem}[main theorem]
For a normed space $P$ the following are equivalent:
\begin{enumerate}
  \item[(i)] $P$ is extremally projective;
  \item[(ii)] $P$ is topologically projective and $\overline P\cong\ellone\Gamma$
    isometrically for some set $\Gamma$;
  \item[(iii)] $\BM(P,\lo\Gamma)=1$ for some set $\Gamma$;
  \item[(iv)] for every $\varepsilon>0$ there is $\Lambda$ with
    $\BM(P,\lo\Lambda)<1+\varepsilon$.
\end{enumerate}
In (iii) necessarily $|\Gamma|=\dim P$, and the distance is realized by
bijections $\overline P\to\ellone\Gamma$ mapping $P$ onto $\lo\Gamma$.
\end{theorem}

\noindent Thus Gr{\o}nb{\ae}k's conjecture holds in the almost isometric form
and fails in the isometric form. For dense subspaces of $\ellone\Lambda$,
topological and extremal projectivity coincide; for uncountable $\Lambda$ this
is a genuine restriction beyond the dimension bound.

\paragraph{Method.} A back-and-forth construction: an automorphism of
$\ellone\Gamma$ of distortion $<1+\varepsilon$ built step by step, moving $P$
exactly onto $\lo\Gamma$. Countable case: one-vector lemma (a block shift of
norm $\le\delta$ swallowing one vector into $\lo{\mathbb N}$), hitting lemma
(a rank-one correction), product lemma. General case: transfinite recursion
along a continuous filtration of $\Gamma$ with perturbations confined to
disjoint blocks, so that they add instead of multiplying; the filtration
comes from closed pairs (elementary-submodel style closing off), and the
blocks are handled by induction on $|\Gamma|$.

\paragraph{Open.} An intrinsic description of the topologically projective
dense subspaces of $\ellone\Lambda$ for uncountable $\Lambda$.
\end{document}
